\documentclass[10pt,a4paper]{article}
\usepackage[hmargin=1.8cm,vmargin=1.5cm,includefoot]{geometry}
\usepackage[T1]{fontenc}
\usepackage[utf8]{inputenc}
\usepackage{lmodern}
\usepackage{microtype}
\usepackage{amsmath}
\usepackage{booktabs}
\usepackage{graphicx}
\usepackage{caption}
\usepackage{subcaption}
\usepackage{xcolor}
\usepackage{tcolorbox}
\usepackage{listings}
\usepackage{enumitem}
\usepackage{placeins}
\usepackage{float}
\usepackage[hidelinks]{hyperref}

\captionsetup{font=small,labelfont=bf,skip=4pt}
\setlist{nosep,leftmargin=*}
\setlength{\parskip}{3pt}
\setlength{\parindent}{0pt}
\makeatletter
\setlength{\@fptop}{0pt}
\setlength{\@fpsep}{12pt}
\setlength{\@fpbot}{0pt plus 1fil}
\makeatother
\newcommand{\figdir}{../output/figures}

\lstset{basicstyle=\ttfamily\scriptsize,breaklines=true,columns=fullflexible,frame=single,
  rulecolor=\color{gray!40},keywordstyle=\color{blue!60!black},commentstyle=\color{gray},
  language=Python,showstringspaces=false}

\title{\vspace{-1.2cm}\textbf{Can we predict stock returns?}\\[2pt]
\large Out-of-sample evidence from U.S. common stocks, 2000-2025}
\author{Ignacio Queipo de Llano P\'erez-Gasc\'on\\\small Predictive Statistics in Finance, IESE Master in Finance}
\date{October 2026}

\begin{document}
\maketitle
\vspace{-0.6cm}

\begin{tcolorbox}[colback=white,colframe=black,boxrule=0.4pt,sharp corners,left=6pt,right=6pt,top=4pt,bottom=4pt]
\textbf{Summary.} No, not in a way that is useful for forecasting returns. Five well-known
stock characteristics have statistically significant links with next month's return in 2001-2019, but together
they explain about 0.13\% of the variation. Frozen at the end of 2019 and tested on 2020-2025, no model beats
forecasting the historical average (1.01\% a month) for every stock. Giving each stock its own average return
improves the historical fit twentyfold and makes the forecasts 5\% worse. The models can still rank stocks a little.
\end{tcolorbox}

\section{The question, the data and the rules of the test}

The firm observes thousands of U.S. stocks at the end of each month and wants to know whether what it sees then
(recent performance, size, risk, trading activity) helps forecast each stock's return over the next month. The data
come from CRSP through WRDS: 1,063,590 stock-month rows for 12,308 common stocks listed on the NYSE, NYSE American or
Nasdaq between January 2000 and December 2025, with a month-end price of at least \$5.

\textit{What one row means.} Each row is one stock in one month. I write $R_{i,t}$ for the total return
(price change plus dividends) of stock $i$ during month $t$; $R_{\text{Apple, Mar 2020}}=-7\%$ is one row. The goal
is to forecast $R_{i,t+1}$, the same stock's return in the following month, using only information available at the
end of month $t$. Stocks are tracked by their permanent CRSP number (PERMNO), because tickers and names change.

\textit{The panel has holes, and that matters.} Stocks list, delist and sometimes leave the sample for a few months
(for example when the price drops below \$5) and come back. There are 14,602 such gaps. If the data were shifted
row by row, the last month before a gap would be paired with a return from several months later and treated as
next month's return: a false target. I therefore place every stock on a calendar grid of stocks by months, so
the previous and the next month always mean calendar months, and a missing month stays missing
(Appendix~\ref{app:data}, Figures~\ref{fig:panel}-\ref{fig:gap}).

\textit{A chronological split.} The models are estimated on 2001-2019 (training) and judged on 2020-2025 (test),
as an investor standing at 31 December 2019 would have to. Each observation is a bet with two dates: the day it is
placed ($t$) and the month it is settled ($t+1$). A bet goes to training only if it was settled by December 2019.
The forecast made in December 2019 for January 2020 therefore belongs to the test sample, because on 31 December
2019 the January return was not yet known. It is one month out of 312, but it keeps any future information out of
training. A random 80/20 split would be less convincing: it would train on 2021 to forecast 2015, which no investor
can do; the model will be used forward in time, so the honest test runs from past to future; and
stocks in the same month move together (in March 2020 almost everything fell): splitting one month between training
and test lets the model learn that month's common shock from half the stocks and predict it for the other half,
which flatters the test. Figure~\ref{fig:split} shows the resulting samples.

\begin{table}[h]
\centering\small
\begin{tabular}{lrrrl}
\toprule
& Stock-months & Stocks & Months & Forecast dates $t$ \\
\midrule
Training (estimation) & 660,763 & 7,823 & 227 & Jan 2001 to Nov 2019 \\
Test (evaluation) & 190,512 & 4,279 & 72 & Dec 2019 to Nov 2025 \\
\bottomrule
\end{tabular}
\caption{Usable samples: rows with a next-month return and all five predictors, the same for every model. Most of
the 212,315 dropped rows are a stock's first year.}
\label{tab:samples}
\end{table}

\section{Five signals, and what I expected before looking}

For each stock and month I build five characteristics, all known at the end of month $t$
(Figure~\ref{fig:windows}):
\begin{itemize}
\item \textbf{Current return} $R_{i,t}$: the return during the month just ended.
\item \textbf{Momentum 12-2}: the compounded return from $t-12$ to $t-2$, requiring all 11 months. The last month is
skipped on purpose: it is the current return, and over one month returns tend to reverse while over a year they tend
to persist. Mixing both in one number would blur two opposite effects.
\item \textbf{Size}: market capitalization (price times shares outstanding).
\item \textbf{Volatility}: the standard deviation of the stock's own last 12 monthly returns (at least 8 observed).
\item \textbf{Turnover}: shares traded in the month divided by shares outstanding, i.e.\ the fraction of the company
that changed hands.
\end{itemize}

\textit{From raw values to monthly ranks.} Each characteristic is converted into its percentile rank among all stocks
that month: 0 is the lowest, 0.5 the median, 1 the highest. The five inputs then share one scale, and one extreme
value (Apple's \$1.3 trillion next to a \$20 million company) cannot dominate the regression
(Figure~\ref{fig:ranks}). Taking logarithms first would change nothing, because a logarithm keeps the order.

The rank also changes what the current-return signal measures, which is why its coefficient is not that of a
standard autoregressive AR(1) model, $R_{i,t+1}=a+bR_{i,t}$. Suppose Apple falls 7\%. If the whole market fell 20\%
that month, Apple was one of the best performers and its rank is close to 1. If the market was flat, the same $-7\%$
puts Apple near the bottom, with a rank close to 0. An AR(1) sees $-0.07$ in both cases. The rank measures how a
stock did relative to other stocks that month; it removes the market-wide move and the exact size of the move (a
fall of 30\% and one of 60\% can both rank near 0).

\begin{table}[h]
\centering\small
\begin{tabular}{p{2.3cm}cp{7.6cm}ccc}
\toprule
Signal & Expected & My reasoning, before looking at the data & Prior & Result & t-stat \\
\midrule
Current return & $-$ & Over one month the rebound dominates the trend: stocks that fell hard tend to recover. & weak & $-$ & $-21.8$ \\
Momentum 12-2 & $+$ & A one-year trend tends to continue. & strong & $\approx 0$ & $-0.8$ \\
Size & $-$ & Small firms are riskier and less liquid, so investors demand a higher return. & weak & $-$ & $-13.4$ \\
Volatility & $-$ & Very volatile stocks attract speculators who overpay for them. & strong & $-$ & $-3.1$ \\
Turnover & $-$ & Rarely traded stocks pay an illiquidity premium, although heavy trading can also signal news. & weak & $\approx 0$ & $1.3$ \\
\bottomrule
\end{tabular}
\caption{The signs I expected, written down before estimating anything so that the explanation could not be
invented afterwards, against the estimates of model M3 on the training sample (Section~\ref{sec:rich}).
A t-statistic beyond $\pm 2$ means the effect is distinguishable from zero at the 95\% level.}
\label{tab:hyp}
\end{table}

\section{One signal: short-term reversal}

The simplest model regresses next month's return on the rank of this month's return, using the training sample only:
$R_{i,t+1}=\alpha+\beta\,\text{Rank}_t(R_{i,t})+u_{i,t+1}$. It is estimated by ordinary least squares (OLS), the line
that minimizes the squared forecast errors. Standard errors use the HC1 correction, which allows the size of the
unexplained noise to differ across observations, as it does between a utility and a biotech start-up.

\begin{center}\small
\begin{tabular}{cccccc}
\toprule
$\hat\beta$ & Std.\ error (HC1) & t-stat & p-value & 95\% interval & In-sample $R^2$ \\
\midrule
$-0.0126$ & 0.00056 & $-22.5$ & $\approx 10^{-112}$ & $[-0.0137,\,-0.0115]$ & 0.098\% \\
\bottomrule
\end{tabular}
\end{center}

\textit{Economic size.} A stock at the 10th percentile of this month's returns is forecast to earn
$0.0163+0.10\times(-0.0126)=1.51\%$ next month; one at the 90th percentile, $0.50\%$. The gap is
$0.8\hat\beta=-1.01$ percentage points a month. The sign means reversal: this month's losers do better than this
month's winners, as I expected. On average the effect is sizeable (about 12\% a year before costs), but it is small
next to the noise, since the typical stock's monthly return has a standard deviation of 11.2\%. Capturing it would
also mean rebuilding the portfolio every month, often in small illiquid stocks, and trading costs would absorb a
large part.

\textit{Significant is not the same as predictive.} The t-statistic is enormous and the $R^2$ is tiny, and the two do
not contradict each other. The p-value answers one narrow question: if the true effect were zero, how likely would
an estimate this far from zero be? Here it is about $10^{-112}$, so the effect is almost certainly real. It
says nothing about whether the effect is large or useful. The t-statistic is $\hat\beta$ divided by its standard
error, and the standard error shrinks as the sample grows; with 660,763 observations even a small effect is measured
very precisely. The $R^2$, the share of the variation in returns that the model explains, does not grow with the
sample. It depends on how much the forecast moves (between 0.5\% and 1.5\%) relative to how much returns move
($\pm$11\%). Collecting more data would not raise it. HC1 treats rows as independent; since stocks in the same month move together, the true t-statistic is somewhat
smaller, but still large.

\section{Adding signals}
\label{sec:rich}

Three nested models are estimated on exactly the same 660,763 rows, so any change in fit comes from the added
signals and not from a different sample: M1 (current return), M2 (adds momentum and size) and M3 (adds volatility
and turnover).

\begin{table}[h]
\centering\footnotesize\setlength{\tabcolsep}{3.5pt}
\begin{tabular}{lccccccc}
\toprule
& Signals & In-sample $R^2$ & Current return & Momentum & Size & Volatility & Turnover \\
\midrule
M1 & 1 & 0.098\% & $-0.0126$ ($-22.5$) & & & & \\
M2 & 3 & 0.132\% & $-0.0122$ ($-21.8$) & $-0.0004$ ($-0.7$) & $-0.0072$ ($-16.0$) & & \\
M3 & 5 & 0.134\% & $-0.0122$ ($-21.8$) & $-0.0005$ ($-0.8$) & $-0.0079$ ($-13.4$) & $-0.0019$ ($-3.1$) & $+0.0009$ ($1.3$) \\
\midrule
\multicolumn{3}{l}{M3: 90th minus 10th percentile, pp/month} & $-0.98$ & $-0.04$ & $-0.63$ & $-0.15$ & $+0.07$ \\
\bottomrule
\end{tabular}
\caption{Pooled OLS on the training sample. Coefficients with HC1 t-statistics in parentheses.}
\label{tab:rich}
\end{table}

Three of my five signs are confirmed (Table~\ref{tab:hyp}). Reversal, which I rated as uncertain, is the strongest
effect. Momentum, my firmest prior, is zero in this period. A likely reason is the momentum crash of 2009: momentum
buys recent winners and sells recent losers, and in the spring of 2009 the biggest losers of the 2008 crisis
(banks, cyclicals) rebounded hardest, wiping out the strategy's gains of the preceding years (Daniel and Moskowitz,
2016). Size is strongly significant ($t=-16$) yet moves the $R^2$ by only 0.03 points, the same lesson as before.

The in-sample $R^2$ rises from 0.098\% to 0.134\%, and it could not have fallen: M1 is M3 with the extra
coefficients set to zero, and OLS picks the best-fitting coefficients, so it always matches the smaller model at
least. It would rise even with a pure-noise variable, which OLS uses to fit this sample's noise, so a higher
in-sample $R^2$ says nothing about forecasting. Only new data can tell the two apart.

\textit{Why the $R^2$ is so low.} Missing variables are only part of the story. First, 16.4\% of the variation in
next-month returns is the move common to all stocks in that month; signals ranked within the month can only say which
stocks will beat the others, never whether the market as a whole rises or falls. Second, most of the rest is news
specific to each company (earnings, lawsuits, takeovers) that is unknowable at $t$. Third, funds would trade away a pattern
that simple. The best published models, with hundreds of signals and neural networks, reach an out-of-sample $R^2$ of
about 0.4\% a month for individual stocks (Gu, Kelly and Xiu, 2020); 0.5\% counts as a success.

\section{The real test: forecasting 2020-2025}

The coefficients of M1-M3 are frozen at the end of 2019 and used to forecast every stock-month of the test period.
Each model is compared with a benchmark that forecasts the same number for every stock: the average next-month
return in training, 1.01\%. (This is exactly the OLS model with only a constant.) Accuracy is measured by the mean
squared error (MSE), the average of the squared forecast errors, and summarized by the out-of-sample $R^2$,
$R^2_{\text{OOS}} = 1-\text{MSE}_{\text{model}}/\text{MSE}_{\text{benchmark}}$. A positive value means the model's forecasts were closer to what happened than the average; a negative one means they
were further away. Unlike the in-sample $R^2$, it can be negative, because the frozen coefficients were chosen to fit
2001-2019 and nothing forces them to beat the average on data they never saw.

\begin{figure}[h]
\centering
\begin{minipage}[t]{0.49\linewidth}\vspace{0pt}\centering\footnotesize\setlength{\tabcolsep}{3pt}
\begin{tabular}{lccc}
\toprule
& In-sample $R^2$ & Test MSE & $R^2_{\text{OOS}}$ \\
\midrule
Benchmark (1.01\%) & n/a & 0.018262 & 0 \\
M1 & 0.098\% & 0.018263 & $-0.002\%$ \\
M2 & 0.132\% & 0.018260 & $+0.010\%$ \\
M3 & 0.134\% & 0.018263 & $-0.004\%$ \\
\bottomrule
\end{tabular}\\[6pt]
\begin{tabular}{lcccccc}
\toprule
$R^2_{\text{OOS}}$ & 2020 & 2021 & 2022 & 2023 & 2024 & 2025 \\
\midrule
M1 & $+.03$ & $+.00$ & $+.05$ & $+.07$ & $-.03$ & $-.17$ \\
M2 & $+.02$ & $-.02$ & $+.14$ & $+.04$ & $-.04$ & $-.13$ \\
M3 & $-.02$ & $-.01$ & $+.16$ & $-.01$ & $-.04$ & $-.15$ \\
\bottomrule
\end{tabular}
\captionof{table}{Test period, 190,512 stock-months. Bottom: $R^2_{\text{OOS}}$ by year, in \%.}
\label{tab:oos}
\end{minipage}\hfill
\begin{minipage}[t]{0.49\linewidth}\vspace{0pt}\centering
\includegraphics[width=\linewidth]{\figdir/f2_r2.pdf}
\caption{Fit keeps rising as signals are added; forecasting accuracy stays at zero.}
\label{fig:r2}
\end{minipage}
\end{figure}

No model beats the average in any meaningful way. The MSEs differ in the sixth decimal, M2 leads by 0.01\% (fifty
times below what the literature calls success), the ranking flips from year to year, and the model with the highest
in-sample $R^2$ (M3) does not have the lowest test error. A negative $R^2_{\text{OOS}}$ means that the model's return
forecasts left us further from reality than assuming every stock earns its usual return.

Figure~\ref{fig:pred} shows why. The forecasts all sit between 0.2\% and 1.8\% while realized returns range from
$-32\%$ to $+41\%$ (1st to 99th percentile); monthly averages across stocks swung between $-21\%$ and $+17\%$ in the
test period, moves that no cross-sectional signal can see. Panel (b) adds a nuance: sorted
into ten groups by forecast, the top group still earned 1.72\% a month against 0.75\% for the bottom group in the test
period (1.61 points in training). The $R^2_{\text{OOS}}$ asks whether the forecast gets the level of each return
right; it does not ask whether the model ranks stocks well. A long-short portfolio could still earn part of that
spread before costs, which is a different question from the one tested here.

\begin{figure}[h]
\centering
\includegraphics[width=0.84\linewidth]{\figdir/f3_pred_vs_real.pdf}
\caption{Forecasts against reality for M3. (a) Distribution of forecasts and realized returns in 2020-2025.
(b) Stocks grouped into deciles by their forecast: average forecast against average realized return.}
\label{fig:pred}
\end{figure}

\begin{figure}[h]
\centering
\includegraphics[width=0.84\linewidth]{\figdir/f6_horizon.pdf}
\caption{The frozen models used month after month, 2020-2025. (a) Growth of \$1 in the average stock against what
each model forecast. (b) Cumulative squared error of the historical mean minus that of each model. (c) Monthly
portfolio long the 10\% of stocks with the highest forecast and short the 10\% with the lowest, before costs.}
\label{fig:horizon}
\end{figure}

Figure~\ref{fig:horizon} follows the frozen models over the six test years. All three forecast a smooth 1\% a month, like the historical mean, while \$1 in the average stock fell to 0.70 in March
2020, rose to 1.45 by 2021 and ended at 1.88 against the 2.05 the models expected (panel a). Panel (b) accumulates the
squared error each model saves relative to the mean: all three were ahead until mid-2023, then gave the advantage
back, and by 2025 M1 and M3 were behind. Used only to rank stocks (panel c), M2 and M3 would have turned \$1 into
about \$1.45 over six years before costs (0.6\% a month on average), with monthly swings of $\pm$4.6\%, a worst month
of $-9\%$ and most of the gain lost after 2024. What works is small, noisy and fragile.

\textit{Three kinds of $R^2$.} Think of an exam. The in-sample $R^2$ is the grade on a practice exam whose questions
you studied, part knowledge and part memorized answers. The population $R^2$ is what you really know; nobody sees it,
and the sample $R^2$ overstates it. The out-of-sample $R^2$ is the grade on the real exam, with new questions and
perhaps a new syllabus. Only it uses unseen data, and only it can reveal a world that has changed.

\section{Why the forecasts failed: the relationships moved}

With 660,763 observations and five signals, classic overfitting is negligible (the memorized part of the
in-sample $R^2$ is about $5/660{,}763\approx0.001\%$), so 0.134\% was a fair estimate for 2001-2019. To see what
changed I re-estimated M3 on the test period, as an after-the-fact diagnostic not used to choose anything
(Figure~\ref{fig:coefs}). Reversal and the size effect survive at
about half their strength. Volatility reverses sign: from $-0.15$ to $+0.83$ points ($t=+7.5$), as highly volatile
stocks led the post-COVID rebound and the speculative rally of 2020-2021. The frozen M3 was betting against the
volatile stocks exactly when they won, which is why it did worse than M2, which does not use volatility. The average
return did not move (1.01\% in both periods), so the problem is the slopes; noise rose (standard deviation 13.5\%
against 11.2\%); and even fitted on 2020-2025 itself, M3 explains only 0.07\%.

\section{Letting each stock have its own average: fixed effects}

Pooled OLS gives every stock the same intercept $\alpha$. A stock fixed-effects model gives each training stock
its own average return, $R_{i,t+1}=\alpha_i+\beta'X_{i,t}+u_{i,t+1}$: 7,823 intercepts. I estimate it by subtracting each stock's own training mean from its returns and signals (the
within transformation), estimating $\beta$, and recovering $\hat\alpha_i=\bar R_i-\hat\beta'\bar X_i$.

\begin{figure}[h]
\centering
\begin{minipage}[t]{0.47\linewidth}\vspace{0pt}\centering
\includegraphics[width=\linewidth]{\figdir/f4_coefs.pdf}
\caption{M3 estimated separately on each period, as the predicted return gap between a 90th and a 10th percentile
stock. Bars show 95\% intervals.}
\label{fig:coefs}
\end{minipage}\hfill
\begin{minipage}[t]{0.47\linewidth}\vspace{0pt}\centering
\includegraphics[width=0.92\linewidth]{\figdir/f5_fe_persistence.pdf}
\caption{Each dot is a stock: its average monthly return in 2001-2019 against 2020-2025 (clipped at $\pm$10\%).}
\label{fig:fe}
\end{minipage}
\end{figure}

\begin{center}\small
\begin{tabular}{lccc}
\toprule
& In-sample $R^2$ (overall) & In-sample $R^2$ (within) & $R^2_{\text{OOS}}$, same 159,605 test rows \\
\midrule
Pooled OLS, M3 & 0.134\% & n/a & $+0.017\%$ \\
Stock fixed effects, M3 & 2.929\% & 1.317\% & $-4.96\%$ \\
\bottomrule
\end{tabular}
\end{center}

The overall $R^2$ compares the full fitted values (own intercept plus signals) with actual returns, so it includes
what the intercepts explain about why some stocks earned more on average. The within $R^2$ removes each stock's mean
first and only measures how well the signals explain each stock's moves around its own average. The overall fit can only improve on pooled OLS: pooled OLS is the special
case with all intercepts equal, and OLS would choose that if it fitted best.

The improvement does not survive. On the test rows of stocks already seen in training, fixed effects are 4.96\% worse
than the average, against $+0.017\%$ for pooled OLS. The reason is Figure~\ref{fig:fe}: across 2,879 stocks, the correlation between a stock's average return in 2001-2019
and in 2020-2025 is only 0.10. Persistent differences in average return across individual stocks are small next to
the noise, and a stock's own history is mostly luck. The intercepts spread 6.7 points a month across stocks, far too much to be
real differences in expected return: half the stocks have under 59 months of history and 28\% under 24, and the
average of 20 returns moving $\pm$11\% says very little. Two further problems: a stock listed after 2019 has no intercept at all (30,907 test rows, 16\%, from about 1,400 new
stocks, many of them the 2021 wave of IPOs and SPACs), and the size coefficient jumps to $-0.152$ ($t=-65$) because
subtracting each stock's mean over the whole period lets the model see, within the sample, whether the stock is
expensive relative to its own future average. That inflates the fit and is useless for forecasting. A workable
alternative gives new stocks the common intercept and pulls each $\alpha_i$ toward it, more so the shorter its
history.

\section{Conclusions for the committee}

\begin{enumerate}
\item The signals are real but small: reversal and size are highly significant yet explain about 0.1\% of returns.
\item Better fit did not mean better forecasts: out of sample every model matched the historical average.
\item The failure came from unstable relationships more than from overfitting: reversal and size halved and volatility
changed sign after 2019. Stock fixed effects add real overfitting and do much worse.
\item I would not use these models to forecast return levels. The ranking signal deserves a proper test: long-short
portfolios net of trading costs, re-estimated as time passes and judged on unseen data.
\end{enumerate}

{\footnotesize\textbf{References.} Ang, Hodrick, Xing and Zhang (2006), \textit{Journal of Finance}; Banz (1981),
\textit{Journal of Financial Economics}; Daniel and Moskowitz (2016), \textit{Journal of Financial Economics}; Datar,
Naik and Radcliffe (1998), \textit{Journal of Financial Markets}; Gu, Kelly and Xiu (2020), \textit{Review of
Financial Studies}; Jegadeesh (1990), \textit{Journal of Finance}; Jegadeesh and Titman (1993), \textit{Journal of
Finance}.\par}

\clearpage
\appendix
\setcounter{figure}{0}\renewcommand{\thefigure}{A\arabic{figure}}
\section{How the data were built, illustrated}
\label{app:data}

\begin{figure}[H]
\centering
\includegraphics[width=0.95\linewidth]{\figdir/f1_split.pdf}
\caption{Usable observations per month. Grey rows are dropped because the stock lacks a next month or a full
momentum history. The black line is the cut between training and test.}
\label{fig:split}
\end{figure}

\begin{figure}[H]
\centering
\includegraphics[width=0.95\linewidth]{\figdir/a1_panel.pdf}
\caption{The unbalanced panel. Each row is a stock, each column a month. White gaps inside a row are months in
which the stock left the sample and later returned.}
\label{fig:panel}
\end{figure}

\begin{figure}[H]
\centering
\includegraphics[width=0.95\linewidth]{\figdir/a2_gap.pdf}
\caption{Building the target for a real stock with a gap (Cypress Semiconductor, 2008-2009). Numbers are actual
monthly returns. Shifting rows pairs October 2008 with the February 2009 return; the calendar grid leaves it without a
target.}
\label{fig:gap}
\end{figure}

\begin{figure}[H]
\centering
\includegraphics[width=0.95\linewidth]{\figdir/a3_windows.pdf}
\caption{Months used by each variable at date $t$. Everything left of the dotted line is known at the end of $t$.
Momentum requires all 11 months; volatility requires 8 of 12.}
\label{fig:windows}
\end{figure}

\begin{figure}[H]
\centering
\includegraphics[width=0.95\linewidth]{\figdir/a4_ranks.pdf}
\caption{Size in December 2019: raw, logarithm and monthly percentile rank. The three marked stocks keep the same
order in all panels.}
\label{fig:ranks}
\end{figure}

\begin{figure}[H]
\centering
\includegraphics[width=0.75\linewidth]{\figdir/a5_funnel.pdf}
\caption{From raw rows to usable observations. Once the 11-month momentum exists, at least 10 of the 12 volatility
months exist too, so volatility removes no further rows. Requiring 8 of 11 momentum months instead would add about
55,000 rows but mix periods of different length.}
\label{fig:funnel}
\end{figure}

\textbf{Data notes.} The file has no missing values in the variables used. One return of $-100\%$ (a SPAC in March
2024, in the test period) was kept: changing test data after looking at it is what the chronological split is meant
to prevent, and its weight among 190,512 observations is negligible. Units of volume and shares outstanding do not
matter, since multiplying by a constant does not change a rank.

\FloatBarrier
\section{Python code}

The Python code that reproduces every number and figure in this report is at
\url{https://github.com/iqueipopg/iese-mif/tree/main/predictive-statistics/stock-returns-case}. The README there
lists the order in which to run the scripts. The CRSP data are licensed through WRDS and are not included.

\end{document}
